\documentclass[10pt]{article}
\usepackage[T1]{fontenc}
\usepackage[margin=0.8in]{geometry}
\pagestyle{empty}
\setlength{\parindent}{0pt}
\setlength{\parskip}{2pt}
\begin{document}

\textbf{Qiong Yan; Peide Yang (corresponding author)}\par
Fujian Jiangxia University, Fuzhou 350108, China; Xiamen University Tan Kah Kee College, Xiamen, China\par
Emails: 381007382@qq.com (Qiong Yan); yangpd@xujc.com (Peide Yang, corresponding author)\par
\vspace{2pt}
09 November 2026\par
\vspace{2pt}
\textbf{To the Editor-in-Chief, Solar Energy (Elsevier)}\par
\vspace{2pt}
Dear Editor,\par
\vspace{2pt}
Please consider the enclosed original research article ``A differentiable, Poynting-consistent transfer-matrix and drift--diffusion kernel for co-design of thin-film solar cells'', by Qiong Yan and Peide Yang, for publication in Solar Energy. It is original work, has not been published previously and is not under consideration elsewhere.\par
\vspace{2pt}
Solar Energy regularly publishes transfer-matrix optical design of thin-film cells and drift-diffusion/SCAPS device optimisation. Our contribution joins these two lines of work: a fully coherent, differentiable transfer-matrix solver is coupled directly to a self-consistent one-dimensional drift-diffusion solver, without a surrogate model and without interpolating the optical field onto the electrical mesh. Existing differentiable photovoltaic solvers with an explicit optical model (dPV, DriftJax) use Beer-Lambert absorption and therefore cannot represent standing-wave interference; surrogate-based digital twins replace the optical stage by a learned model.\par
\vspace{2pt}
The generation rate is Poynting-consistent; the optical model satisfies the absolute energy balance to 1.44e-6 and the Beer-Lambert limit to 1.57e-8. Implicit differentiation of the converged electrical state costs one transposed Jacobian solve; a full gradient evaluation costs 1.10 times a forward solve (about 10\% overhead), and agrees with fourth-order Richardson finite differences across the optical, electrical and composite paths. Three-parameter inverse design recovers all parameters to within 0.0022\% in 44 evaluations.\par
\vspace{2pt}
We would draw attention to the reporting of limitations. The derivative-free benchmark uses true function-call counts; the pre-specified tenfold reduction in eight dimensions is not met (about threefold at moderate accuracy, neither optimiser reaching 1e-4). The production mesh carries a 1.52\% discretisation uncertainty. All data are synthetic: no experimental calibration and no efficiency record are claimed. A SCAPS-1D 3.3.12 comparison was executed and is reported with its pre-specified criteria; two of four criteria pass, and we report it as a limited sanity check, not as validation. We regard this explicit accounting, together with the accompanying reproducible suite, as part of the methodological contribution.\par
\vspace{2pt}
A one-factor 8-D/16-D control (E4-C) and a noisy multi-start comparison against five derivative-free baselines (E5) are reported as independent supplements. They are shipped with their recorded histories and do not revise the negative E4 verdict. An optical-coherence ablation delimits the role of the interference term: coherent and single-pass optics differ by only 0.0071\% in Jsc for the reference device, while the thickness-gradient difference reaches 39.85\% in thin, strongly absorbing stacks. A controlled inverse-design experiment now shows that a single-pass optical model biases the recovered absorber thickness, and a finite-difference-gradient baseline isolates the net benefit of the adjoint; the dense-adjoint cost is measured across five mesh sizes.\par
\vspace{2pt}
Both authors approved the submission. The work is original, has no competing interests and received no specific funding. Generative AI was used only for language editing and manuscript source preparation, as disclosed in the manuscript. The deterministic experiment suite, result files, figure scripts and consistency checker are publicly available from https://github.com/qiantanzhijiao/TMM-DD-Codesign at the time of submission (the same link appears in the manuscript); a frozen Zenodo archive will be linked before publication.\par
\vspace{2pt}
Enclosures: manuscript (LaTeX, PDF, Word), 5 figures, 8 tables, Highlights, source archive with README and the declarations in the manuscript.\par
\vspace{2pt}
Yours sincerely,\par
\vspace{2pt}
\textbf{Peide Yang (corresponding author)}\par
Xiamen University Tan Kah Kee College; yangpd@xujc.com\par

\end{document}
